\documentclass[11pt]{article}

\usepackage[utf8]{inputenc}
\usepackage[T1]{fontenc}
\usepackage[margin=1in]{geometry}
\usepackage{hyperref}
\usepackage{xcolor}
\usepackage{listings}
\usepackage{enumitem}

\hypersetup{
  colorlinks=true,
  linkcolor=blue,
  urlcolor=blue,
  citecolor=blue
}

% Code listing style for copy-pasteable shell commands
\definecolor{codebg}{rgb}{0.96,0.96,0.96}
\lstdefinestyle{shell}{
  basicstyle=\ttfamily\small,
  backgroundcolor=\color{codebg},
  breaklines=true,
  breakatwhitespace=false,
  columns=fullflexible,
  frame=single,
  framesep=4pt,
  showstringspaces=false,
  keepspaces=true,
  xleftmargin=2pt,
  xrightmargin=2pt
}
\lstset{style=shell}

\title{Supplementary User Guide: an end-to-end population-genetics analysis of
1000 Genomes chromosome 22 with VariantFlow}
\author{}
\date{}

\begin{document}
\maketitle

\begin{abstract}
\noindent This supplementary note is a beginner-friendly, end-to-end walkthrough
of a population-genetics analysis of human chromosome~22 from the 1000 Genomes
Project using \textsf{VariantFlow}, a command-line tool for post-calling VCF
operations. It is written for a reader with basic command-line familiarity but no
prior bioinformatics experience. Every command is copy-pasteable; each step
explains what it does and why. Slow steps include a small-region shortcut so the
analysis can be reproduced quickly.
\end{abstract}

\section{Introduction}
A \textbf{VCF} (Variant Call Format) file is the standard text table of genetic
variants: each row is a genomic position where individuals differ, and each column
after the fixed fields is one sample's genotype. \textsf{VariantFlow} performs
\emph{post-calling} VCF operations---the analyses that happen after variants have
been called. Its guiding idea is \textbf{selective execution}: it reads and
computes only the parts of the file a query needs, which keeps filtering and
per-site statistics fast on large files. It reproduces the numbers of established
tools (VCFtools, pixy, scikit-allel) from a single self-contained binary.

By the end of this guide you will have downloaded and indexed a real 1000 Genomes
chromosome-22 VCF, reduced it to clean biallelic SNPs, and computed allele
frequencies, missingness, heterozygosity, Hardy--Weinberg tests, nucleotide
diversity ($\pi$), Tajima's~D, linkage disequilibrium, and Fst between two human
populations---and read what those numbers mean, cautiously. All commands are run
as \texttt{variantflow <command>}.

\paragraph{Time and disk warning.} The full chromosome-22 file is several
gigabytes and some computations take a while. Wherever a step is slow, this guide
offers a small-region shortcut (a ${\sim}1$~Mb slice) so you can follow along in
minutes. Do the small-region version first; scale up once comfortable.

\section{Installation}

\subsection{Build VariantFlow from source}
\textsf{VariantFlow} is written in Rust. You need the Rust toolchain
(\texttt{cargo}); if you do not have it:
\begin{lstlisting}[language=bash]
# Install the Rust toolchain (rustup); accept defaults.
curl --proto '=https' --tlsv1.2 -sSf https://sh.rustup.rs | sh
source "$HOME/.cargo/env"
rustc --version
cargo --version
\end{lstlisting}

Build the release binary from the repository root:
\begin{lstlisting}[language=bash]
# Standard build (pure-Rust VCF/BCF text path)
cargo build --release
./target/release/variantflow --version
\end{lstlisting}

For BCF input or region-based queries backed by HTSlib (e.g.\
\texttt{--region chr22:1-1000000} on a \texttt{.bcf}), build with the HTSlib
feature:
\begin{lstlisting}[language=bash]
cargo build --release --features htslib-static
\end{lstlisting}

Confirm the tool and list the subcommands:
\begin{lstlisting}[language=bash]
./target/release/variantflow --help
\end{lstlisting}
You should see: \texttt{filter}, \texttt{stats}, \texttt{freq},
\texttt{missingness}, \texttt{hardy}, \texttt{het}, \texttt{fst}, \texttt{pi},
\texttt{pixy}, \texttt{tajima-d}, \texttt{ld}, \texttt{index}, \texttt{diff},
\texttt{convert}. For readability, the rest of this guide writes
\texttt{variantflow}; either add the binary to \texttt{PATH}
(\texttt{export PATH="\$PWD/target/release:\$PATH"}) or type the full path.

\subsection{Companion tools and dependencies}
A few standard tools do the jobs VariantFlow deliberately does not (downloading,
indexing, and heavy VCF surgery). The easiest route is Bioconda:
\begin{lstlisting}[language=bash]
# One-time channel setup (order matters)
conda config --add channels bioconda
conda config --add channels conda-forge

# Companions in a fresh environment
conda create -n popgen -c bioconda -c conda-forge \
  bcftools htslib vcftools wget
conda activate popgen
\end{lstlisting}
Here \texttt{bcftools}/\texttt{htslib} provide \texttt{tabix} and \texttt{bgzip}
(indexing, subsetting, slicing), \texttt{vcftools} enables optional cross-checks,
and \texttt{wget} downloads data. Version-check everything:
\begin{lstlisting}[language=bash]
bcftools --version
tabix --version
bgzip  --version
vcftools --version
wget --version | head -n1
\end{lstlisting}

\section{Getting the data}
We use the 1000 Genomes Project 30$\times$ high-coverage phased panel (3202
samples) aligned to GRCh38, hosted by IGSR/EBI on EBI's public FTP. This file is
\textbf{large} (several GB); Section~\ref{sec:slice} shows a 1~Mb slice for quick
work.

\subsection{Download the full chromosome-22 panel}
\begin{lstlisting}[language=bash]
mkdir -p data && cd data

BASE="https://ftp.1000genomes.ebi.ac.uk/vol1/ftp/data_collections/1000G_2504_high_coverage/working/20220422_3202_phased_SNV_INDEL_SV"
FILE="1kGP_high_coverage_Illumina.chr22.filtered.SNV_INDEL_SV_phased_panel.vcf.gz"

# VCF and its tabix index (.tbi); -c resumes partial downloads
wget -c "${BASE}/${FILE}"
wget -c "${BASE}/${FILE}.tbi"
\end{lstlisting}
If the \texttt{.tbi} is absent, build it and check the file opens:
\begin{lstlisting}[language=bash]
tabix -p vcf "${FILE}"
variantflow stats "${FILE}"
\end{lstlisting}

\subsection{Small-region shortcut (recommended first)}
\label{sec:slice}
Slicing a 1~Mb window with \texttt{tabix} yields a small file that runs every
downstream step in seconds (needs the \texttt{.tbi} index):
\begin{lstlisting}[language=bash]
# Extract chr22:20,000,000-21,000,000
tabix -h "${FILE}" chr22:20000000-21000000 | bgzip > chr22.slice.vcf.gz
tabix -p vcf chr22.slice.vcf.gz
variantflow stats chr22.slice.vcf.gz
\end{lstlisting}

\paragraph{Chromosome naming.} These files use \texttt{chr22} (with the
\texttt{chr} prefix). If a region query returns nothing, check whether your file
uses \texttt{22} instead and adjust (e.g.\ \texttt{22:20000000-21000000}).
Set one variable so you can switch between the slice and the full file:
\begin{lstlisting}[language=bash]
# Follow along fast with the slice:
VCF=chr22.slice.vcf.gz
# ...or run the full analysis:
# VCF="1kGP_high_coverage_Illumina.chr22.filtered.SNV_INDEL_SV_phased_panel.vcf.gz"
\end{lstlisting}

\section{Preprocessing}
Raw call sets contain multiallelic sites, indels, and structural variants; most
classic statistics assume \textbf{biallelic SNPs}.

\subsection{Reduce to biallelic SNPs (bcftools)}
\begin{lstlisting}[language=bash]
bcftools view -m2 -M2 -v snps "$VCF" -Oz -o chr22.snps.vcf.gz
tabix -p vcf chr22.snps.vcf.gz
\end{lstlisting}
Here \texttt{-m2 -M2} keeps sites with exactly two alleles (biallelic),
\texttt{-v snps} keeps SNPs only, \texttt{-Oz} writes bgzip output, and
\texttt{tabix} indexes it.

\subsection{Optional: select a subset of samples}
Per-sample commands accept a plain-text file, \textbf{one sample ID per line}
(blank lines and \texttt{\#} comments ignored), via \texttt{--keep} (use only
these) or \texttt{--remove} (drop these):
\begin{lstlisting}[language=bash]
cat > keep.txt <<'EOF'
# one sample ID per line
HG00096
HG00097
HG00099
EOF
\end{lstlisting}

\subsection{Quality filtering with \texttt{filter}}
\texttt{variantflow filter} selects records with a predicate passed to
\texttt{--where}. The expression language understands \texttt{CHROM},
\texttt{POS}, \texttt{QUAL}, \texttt{FILTER}, \texttt{INFO/<TAG>} (with array
indexing like \texttt{INFO/AF[1]}), and \texttt{FORMAT/<TAG>}, plus the aggregate
\texttt{N\_PASS(...)}.
\begin{lstlisting}[language=bash]
# Keep only PASS sites with a decent quality score
variantflow filter chr22.snps.vcf.gz \
  --where 'QUAL > 30 && FILTER == "PASS"' \
  -o chr22.clean.vcf.gz
tabix -p vcf chr22.clean.vcf.gz

# Common variants only (alt-allele frequency above 5%)
variantflow filter chr22.snps.vcf.gz \
  --where 'INFO/AF[0] > 0.05' -o chr22.common.vcf.gz

# Restrict to a region while filtering (needs an index)
variantflow filter chr22.snps.vcf.gz \
  --region chr22:20000000-21000000 --where 'QUAL > 30' -o chr22.region.vcf.gz

# Require >=2 samples with a well-covered alternate allele
variantflow filter chr22.snps.vcf.gz \
  --where 'N_PASS(FORMAT/AD[1] > 10) >= 2' -o chr22.cohort.vcf.gz
\end{lstlisting}

\paragraph{Note on this dataset.} The 1000 Genomes high-coverage panel is already
a filtered, phased release, so most sites are \texttt{PASS} and \texttt{QUAL} is
not always informative. The filter step matters more on your own freshly called
data; here it mainly demonstrates the syntax. The important cleanup is the
biallelic-SNP reduction. From here on, set:
\begin{lstlisting}[language=bash]
SNPS=chr22.snps.vcf.gz
\end{lstlisting}

\section{Population-genetics calculations}
Each subsection gives \emph{what it measures}, \emph{the command}, and \emph{how
to read the output}. Outputs are tab-separated tables.

\subsection{Allele frequency (\texttt{freq})}
\textbf{What it measures.} How common each allele is at every site---the building
block for almost everything else.
\begin{lstlisting}[language=bash]
variantflow freq "$SNPS" -o chr22.frq
\end{lstlisting}
Columns: \texttt{CHROM POS N\_ALLELES N\_CHR \{ALLELE:FREQ\}}. \texttt{N\_CHR} is
the number of chromosomes with data (twice the number of genotyped diploid
samples). Restrict to a group with \texttt{--keep} / \texttt{--remove}.

\subsection{Missingness (\texttt{missingness})}
\textbf{What it measures.} How much genotype data is absent, per site and per
individual. High missingness biases every other statistic.
\begin{lstlisting}[language=bash]
variantflow missingness "$SNPS" -o chr22.miss
\end{lstlisting}
This writes \texttt{chr22.miss.lmiss} (per-site) and \texttt{chr22.miss.imiss}
(per-individual). \texttt{F\_MISS} is the fraction missing (0 = complete). On the
phased panel it is essentially 0---a useful sanity check.

\subsection{Individual heterozygosity (\texttt{het})}
\textbf{What it measures.} Per-sample observed vs.\ expected homozygosity,
summarised as an inbreeding coefficient $F$.
\begin{lstlisting}[language=bash]
variantflow het "$SNPS" -o chr22.het
\end{lstlisting}
Columns: \texttt{INDV O\_HOM E\_HOM N\_SITES F}. $F\approx 0$ matches
Hardy--Weinberg expectation; $F>0$ means fewer heterozygotes than expected
(possible inbreeding or pooled structure); $F<0$ means more (possible
contamination or mixing groups).

\subsection{Hardy--Weinberg equilibrium (\texttt{hardy})}
\textbf{What it measures.} Whether observed genotype counts match Hardy--Weinberg
proportions, with a chi-square test.
\begin{lstlisting}[language=bash]
variantflow hardy "$SNPS" -o chr22.hwe
\end{lstlisting}
Columns include \texttt{OBS\_HOM\_REF OBS\_HET OBS\_HOM\_ALT} and
\texttt{CHISQ\_HWE}. In a pooled multi-population sample, apparent departures are
often \emph{population structure} (the Wahlund effect), not genotyping error.

\subsection{Nucleotide diversity $\pi$ (\texttt{pi}, windowed)}
\textbf{What it measures.} Average nucleotide differences per site between two
random sequences---a core measure of within-population diversity, reported in
windows.
\begin{lstlisting}[language=bash]
variantflow pi "$SNPS" --window-size 100000 -o chr22.windowed.pi
\end{lstlisting}
Higher $\pi$ means a more diverse region. Use \texttt{--keep} to compute $\pi$
within a single population; $\pi$ mixed across populations is inflated by
between-group differences.

\subsection{Tajima's D (\texttt{tajima-d})}
\textbf{What it measures.} Compares two diversity estimates to detect departures
from neutral, constant-size evolution, per window.
\begin{lstlisting}[language=bash]
variantflow tajima-d "$SNPS" --window-size 100000 -o chr22.tajimaD.tsv
\end{lstlisting}
$D\approx 0$ is consistent with neutrality; $D<0$ (excess rare variants) with
expansion or selection; $D>0$ (excess intermediate-frequency variants) with
balancing selection or a bottleneck. A value alone never proves selection.

\subsection{Linkage disequilibrium (\texttt{ld})}
\textbf{What it measures.} Non-random association between alleles at nearby
sites, as a function of distance.
\begin{lstlisting}[language=bash]
variantflow ld "$SNPS" --max-distance 50000 -o chr22.geno.ld
\end{lstlisting}
LD generally decays with distance. Keep \texttt{--max-distance} modest: the number
of site pairs (and runtime) grows quickly.

\subsection{Fst between two populations (\texttt{fst})}
\textbf{What it measures.} Genetic differentiation between two populations: 0 =
identical allele frequencies, higher = more differentiated.

First build two population files from the 1000 Genomes sample panel, which maps
each sample to its population:
\begin{lstlisting}[language=bash]
PED="https://ftp.1000genomes.ebi.ac.uk/vol1/ftp/data_collections/1000G_2504_high_coverage/20130606_g1k_3202_samples_ped_population.txt"
wget -c "$PED" -O samples.panel
head -n1 samples.panel   # inspect columns first!

# Extract two population lists (one sample ID per line).
# Adjust the awk field number to match the sample-ID column.
awk 'NR>1 && $0 ~ /GBR/ {print $2}' samples.panel > GBR.txt   # British
awk 'NR>1 && $0 ~ /YRI/ {print $2}' samples.panel > YRI.txt   # Yoruba
wc -l GBR.txt YRI.txt
\end{lstlisting}
Different panel files order columns differently; inspect the header and set the
\texttt{awk} fields so each output is a plain list of sample IDs, one per line.
Then compute Fst (Hudson estimator by default; \texttt{--estimator weir-cockerham}
is also available):
\begin{lstlisting}[language=bash]
variantflow fst "$SNPS" --pop GBR.txt --pop YRI.txt -o chr22.GBR_YRI.fst
\end{lstlisting}
Output columns: \texttt{CHROM POS HUDSON\_FST} (per site). Fst${}\approx 0$ means
near-identical frequencies; human continental comparisons are often ${\sim}0.1$;
high-Fst sites are selection candidates, not conclusions. Sites monomorphic in
both groups return \texttt{nan} (undefined)---expected, not an error.

\subsection{Missing-data-aware $\pi$ and dxy (\texttt{pixy})}
\textbf{What it measures.} Windowed within-population diversity ($\pi$) and
between-population divergence (dxy) using the pixy estimator, which stays
unbiased under missing data by counting only non-missing pairwise comparisons.

\textbf{Key requirement:} \texttt{pixy} needs an \textbf{all-sites VCF}---one that
includes invariant (monomorphic) sites, not just variant sites---because correct
$\pi$/dxy is the ratio of differences to comparisons over all callable sites. The
standard 1000 Genomes panel is variant-only, so this command targets an all-sites
call set (e.g.\ produced with \texttt{bcftools mpileup}/\texttt{call} emitting all
sites).
\begin{lstlisting}[language=bash]
# populations file: two whitespace/tab separated columns per line:
#   SampleID    PopulationName
#   HG00096   GBR
#   NA18486   YRI
variantflow pixy allsites.vcf.gz \
  --populations pops.txt \
  --window-size 10000 \
  --out-pi   chr22.pixy_pi.tsv \
  --out-dxy  chr22.pixy_dxy.tsv
\end{lstlisting}
Note that \texttt{pixy} takes the input as a \textbf{positional} argument (not
\texttt{--input}), and writes outputs to \texttt{--out-pi} and \texttt{--out-dxy}
(there is no \texttt{-o}). Read \texttt{chr22.pixy\_pi.tsv} (per-population $\pi$
per window) and \texttt{chr22.pixy\_dxy.tsv} (per-pair dxy per window) as in the
$\pi$ and Fst sections, with missing data handled correctly.

\subsection{Site-frequency spectrum (\texttt{sfs}) --- new / unverified}
\paragraph{Verification note.} The \texttt{sfs} subcommand is present in the
current source tree but \textbf{not} in the prebuilt release binary used to verify
this guide, so the commands below could not be run as-is. Rebuild from source
(\texttt{cargo build --release}) to obtain it; treat the exact flags as
provisional.

\textbf{What it measures.} The site-frequency spectrum (SFS/AFS): a histogram of
how many variant sites have each allele count---the raw signal behind Tajima's~D
and demographic inference.
\begin{lstlisting}[language=bash]
# Unfolded (derived-allele) spectrum
variantflow sfs "$SNPS" -o chr22.sfs.tsv

# Folded (minor-allele) spectrum, when the ancestral allele is unknown
variantflow sfs "$SNPS" --folded -o chr22.sfs_folded.tsv
\end{lstlisting}
Use \texttt{--folded} when alleles cannot be polarised into ancestral/derived. The
unfolded output matches \texttt{scikit-allel}'s \texttt{allel.sfs}; the folded
output matches \texttt{allel.sfs\_folded}.

\section{Interpretation}
A cautious reading of what typical values suggest---none is proof on its own.
\begin{itemize}[leftmargin=*]
  \item \textbf{High $\pi$} means high local diversity. African samples (e.g.\
  YRI) typically show higher $\pi$ than non-African samples, reflecting the older,
  larger ancestral African population and the out-of-Africa bottleneck.
  \item \textbf{Negative Tajima's D} genome-wide is common in humans and
  consistent with population expansion; localised strongly negative D can hint at
  a recent sweep, but confirm independently.
  \item \textbf{Positive Tajima's D} suggests an excess of intermediate-frequency
  variants (balancing selection or a bottleneck).
  \item \textbf{Fst} between human continental groups is usually modest (often
  ${\sim}0.1$): most human variation is within rather than between populations.
  \item \textbf{LD decaying with distance} is the expected baseline; extended LD
  can flag low recombination or selection.
  \item \textbf{Missingness near zero} and \textbf{$F$ near zero} on this curated
  panel are healthy sanity checks; large deviations on your own data usually mean
  a data-quality problem to fix before interpreting anything else.
\end{itemize}

\section{Conclusion}
Starting from a raw 1000 Genomes chromosome-22 VCF, you downloaded and indexed the
data, reduced it to clean biallelic SNPs, and computed a full suite of
population-genetics statistics---allele frequency, missingness, heterozygosity,
Hardy--Weinberg, nucleotide diversity, Tajima's~D, linkage disequilibrium, Fst
between two human populations, and (on an all-sites file) the missing-data-aware
pixy estimates---with a single tool and standard companions. The results reproduce
those of established tools and point toward the classic story of human genetic
variation: high diversity, modest between-population differentiation, and
signatures of past demographic change. Use them as a starting point for careful,
hypothesis-driven analysis rather than as final claims.

\begin{thebibliography}{9}
\bibitem{byrska2022} Byrska-Bishop M, Evani US, Zhao X, et al.
  High-coverage whole-genome sequencing of the expanded 1000 Genomes Project
  cohort including 602 trios. \textit{Cell} 185(18):3426--3440 (2022).
  doi:10.1016/j.cell.2022.08.004
\bibitem{vcftools2011} Danecek P, Auton A, Abecasis G, et al.
  The variant call format and VCFtools. \textit{Bioinformatics}
  27(15):2156--2158 (2011). doi:10.1093/bioinformatics/btr330
\bibitem{bcftools2021} Danecek P, Bonfield JK, Liddle J, et al.
  Twelve years of SAMtools and BCFtools. \textit{GigaScience} 10(2):giab008
  (2021). doi:10.1093/gigascience/giab008
\bibitem{pixy2021} Korunes KL, Samuk K. pixy: Unbiased estimation of nucleotide
  diversity and divergence in the presence of missing data.
  \textit{Molecular Ecology Resources} 21(4):1359--1368 (2021).
  doi:10.1111/1755-0998.13326
\bibitem{scikitallel} Miles A, et al. scikit-allel: Explore and analyse genetic
  variation. Zenodo software package.
  \url{https://github.com/cggh/scikit-allel}
\end{thebibliography}

\end{document}
